\section{Kết quả đạt được}

\subsection{Kết quả tối ưu YOLO}

Đồ án đã xây dựng pipeline có thể tái lập gồm baseline, structured pruning,
fine-tuning, QAT, export ONNX và TensorRT. Các cấu hình được đánh giá theo cùng
dataset và bộ chỉ số, giúp lựa chọn model dựa trên trade-off thay vì chỉ theo FPS.
Kết quả cuối cần bổ sung: cấu hình được chọn \cbs{tên cấu hình}, mức thay đổi
$mAP@50{:}95$ \cbs{giá trị}, speedup \cbs{giá trị} và mức giảm kích thước
\cbs{giá trị} so với baseline.

\subsection{Kết quả Edge Runtime}

Edge Runtime đã tổ chức pipeline camera, preprocessing, TensorRT inference,
postprocessing và result processing với metric theo từng stage. Runtime hỗ trợ
load engine theo manifest, warm-up, health check và xuất telemetry. Kết quả cần
bổ sung gồm FPS, latency đầu cuối, frame drop, GPU memory, power và độ ổn định
trong bài chạy kéo dài trên từng thiết bị.

\subsection{Kết quả Edge Platform}

Nền tảng đã thiết kế và hiện thực các miền Fleet, Model Registry, Deployment,
Command, Incident, Knowledge và Ops; đồng thời xây dựng Edge Agent quản lý
desired state, actual state, model lifecycle và runtime lifecycle. Kịch bản triển
khai, canary, monitoring, incident và rollback tạo thành vòng đời vận hành thống
nhất. Các chỉ số success rate, command latency, deployment time và rollback time
cần được tổng hợp từ Chương 6.

\subsection{Kết quả AI Copilot}

AI Copilot kết hợp Knowledge Retrieval với các công cụ Fleet, Runtime,
Deployment và Incident. Thiết kế tách read-only tool khỏi hành động thay đổi
trạng thái, yêu cầu lập kế hoạch, kiểm tra policy, phê duyệt và audit. Kết luận cuối
cần nêu successful task rate, retrieval accuracy, latency, số tool call trung bình
và kết quả kiểm tra vi phạm quyền.